\section{Discussion}
\label{sec:discussion}

\subsection{Governance as Infrastructure}

The central finding is that \kdd~governance is most effective when it is
implemented as infrastructure rather than convention. Enforcing schema
validation at generation time (\sdd), routing all tool calls through an
audited gateway (\mcp), and requiring evidence citations in every recommendation
produces measurable improvements in traceability, audit coverage, and output
groundedness with negligible overhead in normal operation.

\subsection{The Value of the Digital Twin}

The $23\%$ block rate observed in E07 indicates that a non-trivial fraction of
raw copilot recommendations would have been detrimental if applied without
validation. The twin's ability to provide safer alternatives in $91\%$ of blocked
cases suggests that the system degrades gracefully rather than failing silently.

\subsection{Skill Registries and Consistency}

Experiment E04 confirms that registering validated agent capabilities as reusable
skills reduces output variance across runs. This mirrors findings from software
component reuse research and suggests that skill registries are a practical
mechanism for managing LLM output consistency in production agentic systems.

\subsection{Limitations and Future Work}

The current architecture relies on synchronous telemetry processing. Real-time
deployment at MotoGP bandwidth ($>$500\,Hz per channel) would require an
asynchronous streaming pipeline and sub-200\,ms recommendation latency. Future
work will integrate a streaming layer and benchmark end-to-end latency under
realistic load. Additionally, the digital twin models are currently calibrated
on synthetic data; validation against real (anonymised) session data is a
priority for the next research phase.
